\chapter{Hyperbolic, Elliptic, and Parabolic Regimes in the RSVP Field Equations}
\label{ch:hyperbolic-parabolic-elliptic}

This chapter presents a unified analytic, geometric, and semantic
classification of the RSVP field equations
\[
\delta \mathcal{L}/\delta \Phi = 0, 
\qquad
\delta \mathcal{L}/\delta \mathbf{v} = 0,
\qquad
\delta \mathcal{L}/\delta S = 0,
\]
introduced in Chapters~\ref{ch:rsvp-lagrangian}–\ref{ch:stability}.
Our goal is to demonstrate that the RSVP system is
\emph{intrinsically multi-regime}, exhibiting
hyperbolic, elliptic, and parabolic phases
depending on the local semantic curvature, entropy gradient, and
field–metric coupling.

The resulting analytic structure mirrors that of mixed-type PDEs in
fluid dynamics and general relativity, but with novel features induced by
the semantic manifold structure and entropy production.
The chapter culminates in the
\textbf{RSVP Analytic Phase Transition Theorem},
which states that the PDE type transitions
occur precisely when the \emph{semantic scalar curvature}
crosses canonical critical thresholds.
These phase transitions govern the propagation, diffusion, and constraint
enforcement that drive semantic relaxation and ultimately semantic galaxy formation.

% ============================================================
\section{The RSVP Field System and Principal Symbol}
% ============================================================

Let $\mathcal{M}$ be the semantic manifold
defined in Chapter~\ref{ch:semantic-manifold},
equipped with Riemannian metric $g$, volume form $\mathrm{dVol}_g$,
and RSVP fields $(\Phi, \mathbf{v}, S)$.
The Lagrangian density introduced in Chapter~\ref{ch:rsvp-lagrangian}
is
\[
\mathcal{L}
=
\frac12 |\nabla \Phi|^2
+ \frac{1}{2} \|\nabla \mathbf{v}\|^2
+ \frac{\alpha}{2} |\nabla S|^2
+ \beta \Phi \nabla \cdot \mathbf{v}
+ \gamma S (\nabla \cdot \mathbf{v})
+ \frac{\mu}{2} S^2
- V(\Phi).
\]

The RSVP Euler–Lagrange equations are a coupled second-order PDE system
\[
\mathcal{E}(\Phi, \mathbf{v}, S) = 0.
\]
To classify the system, we compute the principal symbol
$\sigma_{\mathcal{E}}(x,\xi)$.
For a unit covector $\xi_a$, define the directional derivatives
\[
D_\xi \Phi = \xi^a \nabla_a \Phi,
\qquad
D_\xi \mathbf{v}_i = \xi^a \nabla_a v_i,
\qquad
D_\xi S = \xi^a \nabla_a S.
\]
The principal symbol matrix is block-structured:
\[
\sigma_{\mathcal{E}}(x,\xi)
=
\begin{pmatrix}
|{\xi}|^2                & \beta \,\xi \otimes & 0 \\
\beta \,\xi \otimes      & |\xi|^2 \mathbf{I}  & \gamma\, \xi\otimes \\
0                       & \gamma\, \xi\otimes & \alpha |\xi|^2
\end{pmatrix}.
\]

We now analyze the eigenvalues of this symbol.

% ============================================================
\section{Eigenvalue Structure and PDE-Type Classification}
% ============================================================

The RSVP system belongs to the class of \emph{mixed-type} PDEs.
The principal symbol has three eigenvalue branches:
\[
\lambda_\Phi = |\xi|^2,
\qquad
\lambda_{v} = |\xi|^2 \pm \beta |\xi|,
\qquad
\lambda_S = \alpha |\xi|^2 \pm \gamma |\xi|.
\]

\paragraph{Hyperbolic Regime.}
The system is hyperbolic when
\[
\beta^2,\,\gamma^2 > 4\,\alpha|\xi|^2.
\]
In this regime propagation cones exist and characteristics are real.
This corresponds to:
\begin{itemize}
\item semantic wavefront propagation,
\item transport of agency ($\mathbf{v}$),
\item long-range semantic influence.
\end{itemize}

\paragraph{Parabolic Regime.}
The system is parabolic when the mixed terms dominate:
\[
\beta \text{ small},\qquad \gamma \text{ small}, \qquad \alpha > 0.
\]
This yields diffusion-dominated flow:
\begin{itemize}
\item entropy smoothing,
\item semantic relaxation,
\item dissipation of curvature.
\end{itemize}

\paragraph{Elliptic Regime.}
The system is elliptic when
\[
\beta = \gamma = 0,
\]
or when curvature forces suppress propagation.
This yields:
\begin{itemize}
\item instantaneous constraint enforcement,
\item semantic consistency equations,
\item potential-driven equilibria.
\end{itemize}

These three regimes coexist on the same manifold and may transition as
local curvature evolves.

% ============================================================
\section{Semantic Curvature and Regime Transition}
% ============================================================

Let $\mathcal{R} = R[g]$ denote the scalar curvature of the semantic manifold
and let
\[
\kappa_\mathrm{eff}
=
\beta \,\nabla \cdot \mathbf{v}
+ \gamma S
\]
denote the \emph{semantic curvature contribution}
from RSVP fields.

Define the total semantic curvature
\[
\mathcal{K} = \mathcal{R} + \kappa_\mathrm{eff}.
\]

\begin{definition}[Critical Curvature]
The PDE-type transitions occur on the hypersurfaces where
\[
\mathcal{K} = 0,
\qquad
\mathcal{K} = \mathcal{K}_c := \frac{\beta^2}{4\alpha},
\qquad
\mathcal{K} = \mathcal{K}_s := \frac{\gamma^2}{4\alpha}.
\]
\end{definition}

These surfaces are RSVP analogues of sonic hypersurfaces in fluid dynamics
and null surfaces in relativity.  
Across these surfaces the RSVP equations change type:
\[
\begin{cases}
\mathcal{K} < 0 &\Rightarrow \text{hyperbolic},\\
0 < \mathcal{K} < \mathcal{K}_c &\Rightarrow \text{parabolic},\\
\mathcal{K} > \mathcal{K}_c &\Rightarrow \text{elliptic}.
\end{cases}
\]

Thus \emph{semantic curvature governs analytic phase transitions}.

% ============================================================
\section{Mixed-Type Structure and Semantic Transition Hypersurfaces}
% ============================================================

The RSVP system is generically of \emph{Tricomi type}:
\[
\text{hyperbolic for } \mathcal{K}<0, 
\quad
\text{elliptic for } \mathcal{K}>\mathcal{K}_c,
\]
with a transitional parabolic layer in between.

These surfaces produce:
\begin{itemize}
\item semantic shock formation in high-curvature hyperbolic regions,
\item diffusive semantic wells in parabolic zones,
\item global constraint enforcement in elliptic basins.
\end{itemize}

The analogy to physical systems is profound:
\begin{itemize}
\item hyperbolic = propagation of intent,
\item parabolic = entropic smoothing / forgetting,
\item elliptic = semantic coherence enforcement.
\end{itemize}

% ============================================================
\section{Lyapunov Functionals and Regime-Dependent Decay}
% ============================================================

The master Lyapunov functional from Chapter~\ref{ch:stability}
has regime-specific decay estimates:
\[
\frac{d}{dt} \mathcal{L}(\Phi, \mathbf{v}, S)
\le
\begin{cases}
- c_\mathrm{hyp} \|\nabla(\Phi,\mathbf{v},S)\|^2 &\text{hyperbolic},\\[4pt]
- c_\mathrm{par} \|\Delta(\Phi,\mathbf{v},S)\|^2 &\text{parabolic},\\[4pt]
- c_\mathrm{ell} \|(\Phi,\mathbf{v},S)\|^2 &\text{elliptic}.\\
\end{cases}
\]

In each regime the dissipation mechanism differs,
but monotonic decay is guaranteed everywhere by entropy production.

% ============================================================
\section{The RSVP Analytic Phase Transition Theorem}
% ============================================================

\begin{theorem}[RSVP Analytic Phase Transition]
Let $(\Phi, \mathbf{v}, S)$ be a smooth solution to the RSVP system on the
semantic manifold $(\mathcal{M},g)$.
Then the set of points where the PDE type changes is the union of the
semantic curvature hypersurfaces
\[
\Sigma_0 = \{x: \mathcal{K}(x)=0\},
\qquad
\Sigma_c = \{x: \mathcal{K}(x)=\mathcal{K}_c\}.
\]
Across these hypersurfaces the system transitions between
hyperbolic, parabolic, and elliptic regimes,
and the Lyapunov functional remains continuous and strictly decreasing.
\end{theorem}

\begin{proof}[Sketch]
The principal symbol changes signature exactly when its eigenvalues
cross zero, which occurs when the semantic curvature crosses the critical
thresholds.  
Continuity of the Lyapunov functional follows from convexity of $\mathcal{L}$
and positivity of the entropy-production current.
\end{proof}

This theorem is the analytic heart of the RSVP relaxation mechanism.

% ============================================================
\section{Semantic Interpretation}
% ============================================================

The unified analytic–geometric view yields a direct cognitive–semantic
interpretation:

\begin{itemize}
\item \textbf{Hyperbolic}: semantic signals propagate (intent, inference, narrative motion).  
\item \textbf{Parabolic}: semantic content diffuses, blends, or relaxes (memory smoothing, concept drift).  
\item \textbf{Elliptic}: global coherence constraints become dominant (semantic closure, consistency, comprehension).  
\end{itemize}

Semantic galaxies (Part III) emerge precisely when local hyperbolic waves
enter parabolic wells and terminate in elliptic attractors.

% ============================================================
\section{Conclusion}
% ============================================================

Chapter~\ref{ch:hyperbolic-parabolic-elliptic} unified the analytic,
geometric, and semantic perspectives on the RSVP field equations:
\begin{itemize}
\item classification of hyperbolic, parabolic, and elliptic regimes,
\item curvature-driven analytic phase transitions,
\item transition hypersurfaces in semantic geometry,
\item Lyapunov decay across all regimes,
\item cognitive–semantic interpretation of PDE type,
\item preparation for the coupled RSVP–Polyxan dynamics of Part II and III.
\end{itemize}

This chapter completes the analytic backbone of the RSVP theory and
prepares the way for Polyxan hypergraph coupling and semantic galaxy formation.

